\section{Sharp approximation of an original low-energy band}

Fix a bandwidth $\eta>0$. Let $H=\sum_i h_i$ be an open chain of $L$
sites of local dimension $q\ge2$, where each $h_i$ is a positive
contraction supported on its two adjacent sites. Write
\[
    \lambda_0=\min\sigma(H),\qquad
    G(c)=\mathbf1_{(-\infty,\lambda_0+c]}(H)\mathcal H,\qquad
    r_\eta=\dim G(\eta).
\]
For a deficit $v>0$, put
\[
    b=\min\{v,\eta/2\},\qquad g=\min\{b/6,1\},\qquad
    u=1+\log q+\log(1/g).
\]
The results below concern the actual original bands. They assume neither
an internal gap nor invariance of these bands under a clipped filter.
The fixed-bandwidth and nontrivial-local-dimension conventions are stated
in~\cite{gap:alvv16_low_energy_family_scope}. These are scoped consequences
supporting the approximation conclusion of
\cite[Theorem~thm:al-ld, lines 1096--1105]{Arad2016RigorousRG}.

\begin{lemma}[A fixed initial space and the actual refined images]
    \label{lem:ld_sharp_normalized_independent_scales}
    \lean{TNLean.MPS.AreaLaw.exists_fixed_bandwidth_normalized_low_energy_approximation_of_independent_scales}
    \leanok
    There are $K_b>0$ and $C\ge21$, depending only on $\eta$, with
    the following property for every prescribed interval $J$.
    One regional space $W_0$ can be chosen before the accuracy so that
    it is $3/200$-viable for $G(\eta-b/2)$ and
    \[
        \dim W_0\le r_\eta\exp\!\left[
            \frac{K_b}{b/2}
            \bigl(1+\log q+\log(1+1/(b/2))\bigr)^3\right].
    \]
    For each $0<\delta\le1$, put
    $a=16+\log(1/g)+\log(1/\delta)$.
    An integer $\ell\ge1$ can be chosen before the vector, with
    \[
        \ell\le8/g+(7C)^{1/4}g^{-1/4}a^{3/4}.
    \]
    Every unit $\psi\in G(\eta-v)$ then has a unit approximation
    $\phi$ satisfying
    \[
        |\langle\phi,\psi\rangle|\ge1-\delta,\qquad
        \operatorname{SR}_{J:J^c}(\phi)
        \le\dim W_0\,(q\ell^2)^{44\ell}.
    \]
\end{lemma}
\begin{proof}\leanok
    \uses{lem:ld_sharp_regional_polynomial_family}
    Select the initial viable space at deficit $b/2$ before $\delta$.
    A long interval is moved inward by the actual collar construction;
    the regional polynomial family is taken at buffer $b/6$.
    Transport only the two closeness assertions by the canonical
    original-band identities. Their losses are each $\delta/100$.
    The original initial space is therefore $1/20$-viable for the
    upper clipped target. Its actual polynomial image has viability
    error at most $\delta/5$ for the lower clipped target.
    Returning to the original lower band gives error at most
    $21\delta/50\le\delta/2$.
    Normalize a vector in this same image and use its dimension bound.
    For a short interval, take $\phi=\psi$ and bound its exact cut rank
    by $q^{|J|}$. Boundary and empty intervals are included.
\end{proof}

\begin{lemma}[The sharp logarithmic cut-rank bound]
    \label{lem:ld_sharp_logarithmic_cut_rank}
    \lean{TNLean.MPS.AreaLaw.exists_fixed_bandwidth_interval_low_energy_normalized_approximation}
    \leanok
    There is $K_\eta>0$, chosen before the chain, interval, deficit,
    accuracy and vector, such that the preceding $\phi$ may be chosen with
    \[
        \log\operatorname{SR}_{J:J^c}(\phi)
        \le\log r_\eta+
        K_\eta\left(\frac{u^3}{g}
        +g^{-1/4}a^{3/4}u\right)(1+\log a).
    \]
\end{lemma}
\begin{proof}\leanok
    \uses{lem:ld_sharp_normalized_independent_scales}
    Since $g\le b/6$ and $g\le1$,
    the initial exponent is at most $8K_bu^3/g$.
    The actual polynomial count has logarithm at most
    \[
        44(1+\log(11+\beta))
        \bigl(10/g+\beta g^{-1/4}a^{3/4}\bigr)u(1+\log a),
        \qquad\beta=(7C)^{1/4}.
    \]
    Add twice this logarithm to that of the initial dimension.
    This retains exactly one $\log r_\eta$.
\end{proof}

\begin{lemma}[From one cut approximation to marginal mass]
    \label{lem:regional_overlap_to_missing_mass}
    \lean{TNLean.MPS.AreaLaw.one_sub_regional_kyFanNorm_le_twice_error}
    \leanok
    If a unit vector $\phi$ has cut rank at most $D$ and
    $|\langle\phi,\psi\rangle|\ge1-\epsilon$, where $\epsilon\le1$,
    then the marginal of $\psi$ has missing leading $D$-eigenvalue
    mass at most $2\epsilon$.
\end{lemma}
\begin{proof}\leanok
    The overlap squared is bounded by the leading marginal mass.
    Use $(1-\epsilon)^2\ge1-2\epsilon$.
\end{proof}

\begin{lemma}[One normalized sequential approximation]
    \label{lem:normalized_sequential_prefix_mass}
    \lean{TNLean.MPS.AreaLaw.exists_unit_mps_of_uniform_prefix_mass}
    \leanok
    Let $\psi$ be a unit vector and $0<\delta\le1$.
    If every internal prefix marginal has missing leading
    $D$-eigenvalue mass at most $\delta/[8(L+1)]$, then there is
    one unit open-boundary MPS $\Phi$ of bond dimension at most $D$
    with $|\langle\Phi,\psi\rangle|\ge1-\delta$ and
    $\|\Phi-\psi\|^2\le\delta/2$.
\end{lemma}
\begin{proof}\leanok
    \uses{lem:variable_normalized_sequential}
    Project successively onto the actual leading marginal eigenspaces.
    The total missing mass is at most $\delta/8$.
    The sequential projection estimate gives that squared error for the
    resulting vector, which is nonzero. Normalization increases the
    distance by a factor at most two, so the squared error is at most
    $\delta/2$. The product retains rank at most $D$ at every cut.
\end{proof}

\begin{lemma}[Sharp simultaneous MPS approximation]
    \label{lem:ld_sharp_global_mps}
    \lean{TNLean.MPS.AreaLaw.exists_fixed_bandwidth_low_energy_normalized_isOpenMPS_approximation}
    \leanok
    There is $K_\eta>0$, depending only on the fixed bandwidth,
    with the following property. For $0<\delta\le1$, put
    \[
        A=16+\log(1/g)+\log\frac{16(L+1)}\delta.
    \]
    There is an integer $B\ge1$, selected before the vector, satisfying
    \[
        \log B\le K_\eta
        \left(\frac{u^3}{g}+g^{-1/4}A^{3/4}u\right)(1+\log A).
    \]
    Every unit $\psi\in G(\eta-v)$ has one unit open-boundary MPS
    approximation $\Phi$ of bond dimension at most $r_\eta B$,
    with $|\langle\Phi,\psi\rangle|\ge1-\delta$ and
    $\|\Phi-\psi\|^2\le\delta/2$.
\end{lemma}
\begin{proof}\leanok
    \uses{lem:ld_sharp_logarithmic_cut_rank,
        lem:regional_overlap_to_missing_mass,lem:normalized_sequential_prefix_mass}
    Apply the cut-rank bound at accuracy $\delta/[16(L+1)]$ to each
    prefix. Round the same exponential rank factor upward once,
    independently of the vector. Every prefix then has missing mass
    at most $\delta/[8(L+1)]$. Apply the normalized sequential
    construction. The rounding adds at most one to the logarithm,
    which is absorbed by increasing $K_\eta$.
\end{proof}

For the source range $0<v\le\eta$, the already proved capped-width
estimate gives, with $B_\eta=\max\{12,\eta\}$,
\[
    1/g\le B_\eta/v,\qquad
    \log(1/g)\le\log B_\eta+\log(1+1/v).
\]
Thus the displayed ranks have the source's
$v^{-1/4}\log(1/\delta)^{3/4}$ accuracy dependence, with only
logarithmic factors in the regularized parameters and a bandwidth-dependent
constant. Substituting $v=\mu/\log L$ gives its low-density scale.
The statement is an existence theorem and asserts no running time.
